\documentclass{article}
\usepackage[T1]{fontenc}
\usepackage{lmodern}
\usepackage{graphicx}
\usepackage{amsmath,amssymb}
\usepackage[a4paper,margin=2cm]{geometry}
\usepackage[absolute,overlay]{textpos}
\setlength{\TPHorizModule}{1mm}
\setlength{\TPVertModule}{1mm}
\usepackage[indonesian]{babel}
\usepackage{hyperref}
\setlength{\emergencystretch}{2em}
% unit: Halaman 25

\begin{document}

% === Halaman 25 (llm) ===

\section*{Terminologi di dalam Kriptografi}

\begin{enumerate}
    \item \textbf{Pesan}: data atau informasi yang dapat dibaca dan dimengerti maknanya (baik dipersepsi secara visual maupun audial)\
    Nama lain: \textbf{plainteks} (plaintext), \textit{plain-image}, \textit{plain-video}, \textit{plain-video}
\end{enumerate}

Di dalam kriptografi, data dan informasi disebut pesan (\textit{message})

Rupa pesan: teks, gambar, musik, video, tabel, daftar belanja, gambar 3D, sinyal control, dll

\end{document}
